\section{Supplementary Evaluation: Usability and Interactive Development}
\seclabel{evaluation}

The implementation of \Frexlib{} is still in its early stages, and
offers many opportuntities for further engineering work to extend its
functionality, expressiveness, ergonomics, and efficiency.
%
However, we have already carried out some small experiments to assess
user experience and frexlet developer experience to establish that the
approach is feasible, and to identify further directions.

\subsection{Quantitative evaluation}
\idris{} encourages interactive, type-driven development, so it is
important that the checker is responsive when the user modifies the program.
Following~\citet{nielsen:usability}, our \idris{}
implementation aims for response times under one second, and we treat
a response time of over 10 seconds when type checking a modification
to \Frexlib{} client code as a bug.

\pgfplotstableread[col sep = comma]{csv/1vars_commutative.csv}\OneVarCommutativeData
\pgfplotstableread[col sep = comma]{csv/1vars_noncommutative.csv}\OneVarNoncommutativeData
\pgfplotstableread[col sep = comma]{csv/5vars_commutative.csv}\FiveVarCommutativeData
\pgfplotstableread[col sep = comma]{csv/5vars_noncommutative.csv}\FiveVarNoncommutativeData
\pgfplotstableread[col sep = comma]{csv/10vars_commutative.csv}\TenVarCommutativeData
\pgfplotstableread[col sep = comma]{csv/10vars_noncommutative.csv}\TenVarNoncommutativeData
\pgfplotstableread[col sep = comma]{csv/15vars_commutative.csv}\FifteenVarCommutativeData
\pgfplotstableread[col sep = comma]{csv/15vars_noncommutative.csv}\FifteenVarNoncommutativeData

For typical small equalities that arise incidentally in dependently
typed programs, Frex's performance falls very comfortably
within \citeauthor{nielsen:usability}'s limits.  For example, the
checking time\footnote{We use a dated
AMD FX-8320 machine with 16GB memory, running Idris 2 version
0.5.1-1011cc616 on Debian Linux.} is under 0.1s for terms of size six or below with the commutative solver
and terms of size 14 or below with the non-commutative solver, creating an impression of instantaneous response.

As the term size increases, Frex eventually crosses the one second
interactivity threshold.  \figref{timings} shows how type-checking
times grow with term size and with the number of free variables in a
randomly generated term for the commutative and non-commutative monoid
solvers.
%
As the figure shows, Frex's type-checking time generally remains below
the interactivity threshold up to terms of around size 30, and only
exceeds the 10 second threshold (beyond which users' attention is
lost) for a few terms of size 45 or above.
%
Our experience with Frex development suggests that the anomalously
high checking times for these terms is likely to arise from a
performance bottleneck in \idris{}'s evaluator (\secref{lessons})
and that the ongoing development of \idris{} may eventually eliminate
the problem, bringing the type-checking time for most terms up to size
60 down to a few seconds.
%
\begin{figure}[ht]

\begin{tikzpicture}
  \begin{axis}[only marks,width=0.85\columnwidth,height=5.3cm,
    xlabel={term size (leaves)},
    ylabel={run time (s)},
    ymax = 45,
    xmin = 0,
    xmax = 60,
    xlabel style={font=\small},
    ylabel style={font=\small},
    ymode=log,
    log ticks with fixed point,
    scaled ticks=false,
    tick label style={/pgf/number format/fixed},
    title={\footnotesize commutative solver},
    every axis title/.style={below right,at={(0,1)}},
    legend style={anchor=north west,at={(1.05,1.0)},draw=black,fill=white,align=left}
    ]
    \addplot[blue!50!black,mark=triangle,mark size=1pt] table[x index = {0}, y expr=(\thisrow{time})]{\OneVarCommutativeData};
    \addlegendentry{\scriptsize 1 free var};
    \addplot[red!50!black,,mark=otimes,mark size=1pt] table[x index = {0}, y expr=(\thisrow{time})]{\FiveVarCommutativeData};
    \addlegendentry{\scriptsize 5 free vars};
    \addplot[green!50!black,mark=star,mark size=1pt] table[x index = {0}, y expr=(\thisrow{time})]{\TenVarCommutativeData};
    \addlegendentry{\scriptsize 10 free vars};
    \addplot[yellow!50!black,mark=diamond,mark size=1pt] table[x index = {0}, y expr=(\thisrow{time})]{\FifteenVarCommutativeData};
    \addlegendentry{\scriptsize 15 free vars};

    \draw [dashed,green!50!black] (axis cs:0,0.1) -- node[anchor=east,fill=white,pos=0.98, inner sep=1pt] { \footnotesize instantaneity threshold } (axis cs:60,0.1);
    \draw [dashed,green!50!black] (axis cs:0,1) -- node[anchor=west,fill=white,pos=0.02,inner sep=1pt] { \footnotesize interactivity threshold } (axis cs:60,1);
    \draw [dashed,green!50!black] (axis cs:0,10) -- node[anchor=west,fill=white,pos=0.02, inner sep=1pt] { \footnotesize attention threshold } (axis cs:60,10);
  \end{axis}
\end{tikzpicture}


\begin{tikzpicture}
  \begin{axis}[only marks,width=0.85\columnwidth,height=5.3cm,
    xlabel={term size (leaves)},
    xmin=0,
    xmax=60,
    ymax = 30,
    ymode=log,
    log ticks with fixed point,
    ylabel={run time (s)},
    xlabel style={font=\small},
    ylabel style={font=\small},
    scaled ticks=false,
    tick label style={/pgf/number format/fixed},
    title={\footnotesize non-commutative solver},
    every axis title/.style={below right,at={(0,1)}},
    legend style={anchor=north west,at={(1.05,1.0)},draw=black,fill=white,align=left}
    ]
    \addplot[blue!50!black,mark=triangle,mark size=1pt] table[x index = {0}, y expr=(\thisrow{time})]{\OneVarNoncommutativeData};
    \addlegendentry{\scriptsize 1 free var};
    \addplot[red!50!black,mark=otimes,mark size=1pt] table[x index = {0}, y expr=(\thisrow{time})]{\FiveVarNoncommutativeData};
    \addlegendentry{\scriptsize 5 free vars};
    \addplot[green!50!black,mark=star,mark size=1pt] table[x index = {0}, y expr=(\thisrow{time})]{\TenVarNoncommutativeData};
    \addlegendentry{\scriptsize 10 free vars};
    \addplot[yellow!50!black,mark=diamond,mark size=1pt] table[x index = {0}, y expr=(\thisrow{time})]{\FifteenVarNoncommutativeData};
    \addlegendentry{\scriptsize 15 free vars};

    \draw [dashed,green!50!black] (axis cs:0,0.1) -- node[anchor=east,fill=white,pos=0.98, inner sep=1pt] { \footnotesize instantaneity threshold } (axis cs:60,0.1);
    \draw [dashed,green!50!black] (axis cs:0,1) -- node[anchor=west,fill=white,pos=0.02, inner sep=1pt] { \footnotesize interactivity threshold } (axis cs:60,1);
    \draw [dashed,green!50!black] (axis cs:0,10) -- node[anchor=west,fill=white,pos=0.02, inner sep=1pt] { \footnotesize attention threshold } (axis cs:60,10);
  \end{axis}
\end{tikzpicture}

   \caption{\Frexlib{} monoid simplifiers type-checking times}
  \figlabel{timings}
\end{figure}

\subsection{Qualitative evaluation}
To evaluate \Frexlib{}'s usage, we have reproduced
\citeauthor{brady-mckinna-et-al:indexed-binary}'s \citeyearpar{brady-mckinna-et-al:indexed-binary}
dependently typed representation of binary
arithmetic. \citeauthor{brady-mckinna-et-al:indexed-binary} index
binary representations by the natural numbers that they represent, and
so the programmer needs to prove arithmetic operations correct.
Such proofs interleave insightful
equational reasoning steps with rote calculational steps
such as the following:
\begin{center}
\FrexifyLHS{} \IdrisFunction{=} \FrexifyRHS{}
\end{center}
which may be discharged with \IdrisFunction{solve}.
We do not use our reflection capabilities since
these kinds of examples, in which the binding-time analysis is challenging,
are beyond their reach at the moment.
%
With early \Frexlib{} implementations the task was arduous due to
performance bottlenecks in \idris{} that are now eliminated.
%
The only other significant obstacle we encountered was the usual pain
point involved in invoking an algebraic simplifier without a goal
extraction mechanism: the need to repeat the equation and its relevant
rewriting-context when calling \Frexlib{}.
